\begin{textsample}{POS Dim 2 – 2017\_c5 – Score -1.00 – 2017\_c5/enade\_2017\_tec\_an\_sist\_f950c36aa2.txt}  \label{ex:f2_pos_388}
CAPÍTULO 5 DISTRIBUIÇÃO DOS CONCEITOS 5.1 PANORAMA NACIONAL DA DISTRIBUIÇÃO DOS CONCEITOS O Conceito Enade é um indicador de qualidade que avalia os cursos por intermédio dos desempenhos dos estudantes a partir dos resultados do Enade.
Os valores possíveis vão de 1 ( pior situação ) a 5 ( melhor situação ).
Os cursos com apenas um ou sem nenhum concluinte participante não obtêm o Conceito Enade, ficando Sem Conceito ( SC ).
A Tabela 5.1 apresenta a quantidade e distribuição de cursos de Tecnologia em Análise e Desenvolvimento de Sistemas participantes do Enade/2017, por faixa de conceito e Grande Região.
Enfatiza -se, mais uma vez, que a diferença entre os cursos tabulados neste capítulo e no capítulo 2 corresponde aos cursos Sem Conceito ( SC ), ou seja àqueles, em princípio, sem alunos concluintes que houvessem participado da prova ou àqueles que tivessem tão somente um aluno concluinte¹⁸.
Observando -se os dados da Tabela 5.1, nota -se que, dos 386 cursos participantes, 153 ( 39,6 \% ) classificaram -se com conceito 3, o valor modal.
Esse foi também o conceito modal em quase todas as regiões : Nordeste ( 36,2 \% ), Sudeste ( 42,9 \% ), Sul ( 37,4 \% ) e Centro-Oeste ( 42,9 \% ).
Na região Norte, a moda foi o conceito 2, atribuído para 58,8 \% dos cursos da região.
O conceito 2 foi o segundo mais frequente em nível nacional ( 27,2 \%, correspondendo a 105 cursos ) e o conceito 4, o terceiro ( 22,5 \%, correspondendo a 87 cursos ).
Houve, ainda, 12 cursos ( 3,1 \% ) que receberam conceito 5 e dez cursos ( 2,6 \% ) que receberam conceito 1.
Dos 386 cursos de Tecnologia em Análise e Desenvolvimento de Sistemas, 19 ( 4,9 \% ) ficaram Sem Conceito ( SC ). ¹⁸ Estes cursos com somente um concluinte não foram considerados no capítulo 2.
A região Norte participou com 17 cursos ou 4,4 \% do total nacional.
Desses, dez cursos receberam o conceito 2, o que equivale a 58,8 \% do total regional, além de ser o conceito modal na região, como já comentado.
A três cursos ( 17,6 \% ) atribuiu -se conceito 3, a outros três cursos ( 17,6 \% ), conceito 4.
Um curso ficou Sem Conceito ( SC ) e nenhum curso recebeu os conceitos 1 e 5 nessa região.
A região Nordeste participou com 47 cursos ou 12,2 \% do total nacional.
Nessa região, 17 cursos ( 36,2 \% em termos regionais ) obtiveram conceito 3, o conceito modal para a região, como já comentado.
Os conceitos 1 e 2 foram atribuídos, respectivamente, a dois e a 11 cursos ( 4,3 \% e 23,4 \%, respectivamente ).
Ainda nessa região, dez cursos ( 21,3 \% ) foram avaliados com o conceito 4, e quatro cursos ( 8,5 \% ), com conceito 5.
Três cursos ficaram Sem Conceito ( SC ).
Dos 189 cursos participantes da região Sudeste ( 49,0 \% do total nacional ), o conceito 3 foi o conceito modal, como já anteriormente assinalado, tendo sido obtido por 81 cursos ( 42,9 \% ).
O conceito 1 foi atribuído a sete cursos ( 3,7 \% ) e o conceito 2, a 52 cursos ( 27,5 \% ).
Receberam os conceitos 4 e 5, respectivamente, 35 cursos ( 18,5 \% ) e quatro cursos ( 2,1 \% ).
Nessa região, dez cursos ficaram Sem Conceito ( SC ).
Os 91 cursos da região Sul corresponderam a 23,6 \% do total nacional.
Como já indicado, houve predominância do conceito 3, atribuído a 37,4 \% dos cursos da região ( conceito modal ), ou seja, atribuído a 34 dos 91 cursos participantes na região Sul.
O conceito 1 foi atribuído a um curso ( 1,1 \% ), o conceito 2, a 15 cursos ( 16,5 \% ), o conceito 4, a 32 cursos ( 35,2 \% ) e o conceito 5, a quatro cursos ( 4,4 \% ).
Cinco cursos ficaram Sem Conceito ( SC ) na região Sul.
Dos 42 cursos participantes na região Centro-Oeste ( 10,9 \% do total nacional ), 18 receberam conceito 3, o conceito modal ( 42,9 \% ), como já fora destacado.
A 17 cursos ( 40,5 \% ) foi atribuído conceito 2, e a sete cursos ( 16,7 \% ), conceito 4.
Os conceitos restantes, 1 e 5, não foram atribuídos a nenhum curso.
Nessa região, nenhum curso ficou Sem Conceito ( SC ).
O Gráfico 5.1 apresenta a distribuição cumulativa do Conceito Enade segundo Grande Região.
Num gráfico de distribuição cumulativa, em cada valor de abscissa ( eixo x ), o valor apresentado no eixo das ordenadas ( eixo y ) é o valor acumulado da variável.
Por exemplo, a região Norte ( linha azul ) apresenta um pouco mais de 80 \% dos cursos com conceito 3 ou abaixo deste valor.
A região Nordeste ( linha verde ), por outro lado, apresenta um pouco menos de 70 \% dos cursos na mesma situação : com conceito 3 ou abaixo.
Nas regiões Norte ( linha azul ) e Centro-Oeste ( linha preta ), todos os cursos têm conceitos 4 ou abaixo disso, já que a poligonal cumulativa alcança 100 \% para o conceito 4.
Em linhas gerais, num gráfico de distribuição cumulativa, poligonais mais à esquerda representam Grandes Regiões com pior distribuição, e poligonais mais à direita, Grandes Regiões com melhores conceitos.
A região Norte ( linha azul ) apresenta os cursos com conceitos mais baixos, ao passo que as regiões Nordeste ( linha verde ) e Sul ( linha roxa ) apresentam as distribuições com valores maiores.
Cursos em IES situadas nas regiões Sudeste ( linha vermelha ) e Centro-Oeste ( linha preta ) apresentam uma situação intermediária. 5.2 CONCEITOS POR CATEGORIA ADMINISTRATIVA, POR MODALIDADE DE ENSINO E GRANDE REGIÃO A Tabela 5.2 apresenta a distribuição dos cursos de Tecnologia em Análise e Desenvolvimento de Sistemas participantes do Enade/2017, por Categoria Administrativa e Modalidade de ensino, de acordo com os conceitos por eles alcançados, segundo as Grandes Regiões brasileiras.
Dos 386 cursos participantes, 128 ( 33,2 \% ) eram ministrados em Instituições Públicas e 258 ( 66,8 \% ), em Instituições Privadas.
Quanto a Modalidade de ensino, 26 cursos eram a Distância e 360, Presenciais.
De acordo com as informações da Tabela 5.2, em termos nacionais, entre Instituições Públicas e Privadas, observa -se que, dos 12 cursos avaliados com conceito 5, oito eram oferecidos em IES Públicas e quatro, em IES Privadas.
Dos 128 cursos participantes de IES Públicas, o conceito 3 foi o valor modal, atribuído a 54 cursos ( 42,2 \% ).
Entre os demais cursos de Instituições Públicas participantes, um curso foi avaliado com conceito 1, outros 17 cursos ( 13,3 \% cada ) obtiveram conceito 2, 47 cursos ( 36,7 \% ) receberam conceito 4 e, como já comentado, 8 cursos ( 6,3 \% ) receberam conceito 5.
Nesta categoria, um curso ficou Sem Conceito ( SC ).
Na rede Privada, o conceito modal foi 3, atribuído a 99 cursos, o correspondente a 38,4 \% dos 258 cursos da categoria.
Entre os demais cursos participantes, nove ( 3,5 \% ) receberam conceito 1, e 88 ( 34,1 \% ), conceito 2.
O conceito 4 foi atribuído a 40 cursos ( 15,5 \% ) e o conceito 5, a quatro cursos ( 1,6 \% ).
Nesta Categoria Administrativa, 18 cursos ( 7,0 \% ) ficaram Sem Conceito ( SC ).
Dos 360 cursos participantes da modalidade Presencial, o conceito 3 foi o valor modal, atribuído a 141 cursos ( 39,2 \% ).
Entre os demais cursos nesta modalidade de ensino, 84 cursos ( 23,3 \% ) foram avaliados com conceito 4, 96 cursos ( 26,7 \% ) obtiveram conceito 2, 12 cursos ( 3,3 \% ) receberam conceito 5 e dez cursos ( 2,8 \% ) receberam conceito 1.
Nesta modalidade 17 cursos ficaram Sem Conceito ( SC ).
Entre os cursos a Distância, o conceito modal foi 3, atribuído a 12 cursos, o correspondente a 46,2 \% dos 26 cursos da modalidade.
Entre os demais cursos participantes, nove ( 34,6 \% ) receberam conceito 2, três ( 11,5 \% ) receberam conceito 4.
Nessa Modalidade dois cursos ficaram Sem Conceito ( SC ) e nenhum recebeu os conceitos 1 e 5.
Na análise por região, observa -se que, na região Norte, as Instituições Privadas participaram com dez cursos ( 58,8 \% do total regional ), dos quais um ficou Sem Conceito ( SC ) e nenhum recebeu os conceitos 1 e 5.
O conceito modal para as Instituições Privadas na região foi 2, atribuído a seis cursos, correspondendo a 60,0 \% dessa categoria na região.
Dois cursos foram avaliados com conceito 3 e um curso recebeu conceito 4.
As Instituições Públicas participaram com sete cursos na região Norte ( 41,2 \% do total regional ), dos quais quatro obtiveram conceito 2, o valor modal.
O conceito 3 foi atribuído a um curso, e o conceito 4, a dois cursos.
Nenhum curso ficou Sem Conceito ( SC ) ou com conceito 1 ou 5 nessa combinação de categoria e região.
Todos os cursos oferecidos na região Norte eram presenciais.
Na região Nordeste, a rede privada concentrou 33 dos 47 cursos participantes, o equivalente a 70,2 \% do total da região.
Dentre os cursos oferecidos por IES Privadas no Nordeste, 13 cursos foram avaliados com conceito 3, o conceito modal.
Dois cursos obtiveram conceito 1, dez cursos receberam conceito 2, quatro cursos ficaram com conceito 4 e um curso alcançou o conceito máximo, 5.
Nessa combinação de categoria e região, três cursos ficaram Sem Conceito ( SC ).
As Instituições Públicas da região Nordeste participaram com 14 cursos ( 29,8 \% do total da região ).
Desses, seis obtiveram conceito 4, o valor modal, seguidos por outros quatro cursos, que receberam conceito 3.
O conceito 5 foi atribuído a três cursos e o conceito 2 a um curso.
Nenhum curso ficou Sem Conceito ( SC ) ou com conceito 1 nessa região.
A maioria dos cursos na região foram presencias : 44 ( 93,6 \% ).
Para estes, a moda foi o conceito 3 com 15 cursos.
Os demais cursos presenciais da região foram alocados aos conceitos 1 ( dois cursos ), 2 ( 11 cursos ), 4 ( dez cursos ) e 5 ( quatro cursos ).
Dois cursos presenciais ficaram Sem Conceito ( SC ).
A dois dos três cursos a Distância foi atribuído o conceito 3.
Um curso a Distância da região ficou Sem Conceito ( SC ).
Na região Sudeste, a rede privada concentrou 122 dos 189 cursos participantes, o equivalente a 64,6 \% do total da região.
Nessa categoria e região, o conceito modal foi 3, atribuído a 48 cursos.
Nove cursos ficaram Sem Conceito ( SC ).
Os demais foram avaliados com conceito 1 ( seis cursos ), conceito 2 ( 44 cursos ), conceito 4 ( 12 cursos ) e conceito 5 ( três cursos ).
Entre os 67 cursos oferecidos por Instituições Públicas na região Sudeste, o conceito modal foi 3, atribuído a 33 cursos.
Os demais cursos receberam os conceitos 1 e 5 ( ambos com um curso ), 4 ( 23 cursos ) e 2 ( oito cursos ).
Na região Sudeste, nesta categoria, um curso ficou Sem Conceito ( SC ).
Também na região Sudeste, a maioria dos cursos foram presencias : 175 ( 92,6 \% ).
Para estes, a moda foi o conceito 3 com 75 cursos.
Os demais cursos presenciais da região foram alocados aos conceitos 1 ( sete cursos ), 2 ( 46 cursos ), 4 ( 34 cursos ) e 5 ( quatro cursos ).
Nove cursos presenciais ficaram Sem Conceito ( SC ).
Para os 14 cursos a Distância, a distribuição é bimodal nos conceitos 2 e 3 ( 42,9 \% dos cursos em cada conceito ).
Dos outros dois cursos a Distância da região um foi alocado ao conceito 4 e o outro ficou Sem Conceito ( SC ).
Nenhum curso a Distância da região foi alocado aos conceitos 1 ou 5.
As Instituições Privadas concentraram 63 dos 91 cursos participantes da região Sul, 69,2 \% do total regional.
Desses, 26 obtiveram conceito 3, o conceito modal.
Nesta combinação de Categoria Administrativa e Grande Região, os conceitos 2 e 4 foram atribuídos a 14 e a 17 cursos, respectivamente.
Um curso recebeu conceito 1, nenhum recebeu conceito 5 e cinco ficaram Sem Conceito ( SC ).
As Instituições Públicas na região Sul participaram com 28 cursos, 30,8 \% dos cursos da região.
O conceito modal foi 4, atribuído a 15 cursos.
Aos demais cursos foram atribuídos os conceitos 2 ( um curso ), 3 ( oito cursos ) e 5 ( quatro cursos ).
Nenhum curso ficou Sem Conceito ( SC ) ou com conceito 1 nessa região.
Também na região Sul, a maioria dos cursos foram presencias, correspondendo a 85, por oposição aos seis cursos a Distância.
Para os cursos presenciais, a distribuição é bimodal nos conceitos 3 e 4, com 31 cursos cada.
Os demais cursos presenciais da região foram alocados aos conceitos 1 ( um curso ), 2 ( 13 cursos ) e 5 ( quatro cursos ).
Cinco cursos ficaram Sem Conceito ( SC ).
Para os seis cursos a Distância, foram atribuídos os conceitos 2, 3 e 4 com, respectivamente, dois cursos, três cursos e um curso.
Nenhum curso a Distância da região foi alocado aos conceitos 1 e 5 ou ficou Sem Conceito ( SC ).
Na região Centro-Oeste, 30 dos 42 cursos participantes eram de Instituições Privadas ( 71,4 \% em termos regionais ).
Desses, 14 receberam conceito 2, o conceito modal.
Entre os demais, dez receberam conceito 3 e seis cursos receberam conceito 4.
Nenhum curso ficou Sem Conceito ( SC ) ou recebeu os conceitos 1 e 5.
Quanto aos 12 cursos oferecidos por Instituições Públicas na região Centro-Oeste ( 28,6 \% do total regional ), oito cursos ficaram com conceito 3, conceito modal, e três ficaram com o conceito 2.
O conceito 4 foi atribuído a um curso.
Nenhum curso ficou Sem Conceito ( SC ) ou com conceitos 1 ou 5.
Nesta região, 39 cursos foram presencias e três, a Distância.
Para os cursos presenciais, a moda foi o conceito 3 com 17 cursos.
Os demais foram alocados aos conceitos 2 ( 16 cursos ) e 4 ( seis cursos ).
Nenhum curso presencial ficou Sem Conceito ( SC ) ou recebeu os conceitos 1 e 5.
Os três cursos a Distância foram alocados nos conceitos 2, 3 e 4, um curso cada.
O Gráfico 5.2 apresenta a distribuição do Conceito Enade segundo a Categoria Administrativa da IES.
Os cursos em IES Públicas ( linha azul ) apresentam uma poligonal mais à direita do que os em IES Privadas ( linha verde ) e, consequentemente, uma distribuição de Conceitos Enade com valores maiores. 5.3 CONCEITOS POR ORGANIZAÇÃO ACADÊMICA E GRANDE REGIÃO Na Tabela 5.3, encontra -se a distribuição dos conceitos atribuídos aos cursos participantes do Enade/2017 na Área de Tecnologia em Análise e Desenvolvimento de Sistemas, por Organização Acadêmica, segundo as Grandes Regiões brasileiras.
Dos 386 cursos de Tecnologia em Análise e Desenvolvimento de Sistemas participantes, 95 eram oferecidos em Universidades, 67 em Centros Universitários, 150 em Faculdades e 74 em CEFET/IFET.
Esta distribuição corresponde a, respectivamente, 24,6 \%, 17,4 \%, 38,9 \% e 19,2 \% dos cursos.
De acordo com os dados apresentados, de todos os 12 cursos avaliados com conceito 5, um era vinculado a Universidades.
Os cursos vinculados a esse tipo de Organização Acadêmica tiveram o conceito 3 como conceito modal, atribuído a 44 cursos ( 46,3 \% ).
Os demais cursos vinculados a Universidades receberam os conceitos 1 ( dois cursos ), 2 ( 25 cursos ), 4 ( 20 cursos ) e conceito 5 ( um curso, como já mencionado ).
Ficaram Sem Conceito ( SC ), três cursos.
Entre os cursos em Centros Universitários, o conceito modal foi também 3, atribuído a 29 cursos ( 43,3 \% ).
Três cursos vinculados a esse tipo de Organização Acadêmica ficaram Sem Conceito ( SC ), e os outros cursos receberam os conceitos 1 ( dois cursos ), 2 ( 15 cursos ), 4 ( 18 cursos ).
O conceito 5 não foi atribuído a nenhum curso.
Dos 150 cursos mantidos por Faculdades, 57 ( 38,0 \% ) receberam conceito 2, o conceito modal.
O segundo conceito mais frequentemente atribuído aos cursos oferecidos por essa categoria de IES foi o conceito 3, recebido por 49 cursos.
Na sequência, 23 cursos obtiveram conceito 4, cinco cursos receberam o conceito 1 e outros quatro o conceito 5.
Além disso, 12 cursos ficaram Sem Conceito ( SC ).
Dos 74 cursos mantidos por CEFET/IFET, 31 ( 41,9 \% ) receberam o conceito modal 3.
O segundo conceito mais frequentemente atribuído aos cursos oferecidos por essa categoria de IES foi o conceito 4, recebido por 26 cursos.
Foram atribuídos os conceitos 5 para sete cursos e 2 para oito cursos.
Um curso ficou Sem Conceito ( SC ) e um recebeu conceito 1.
Considerando -se separadamente as regiões brasileiras, verifica -se que, na região Norte, as Universidades participaram com dois dos 17 cursos participantes.
O conceito modal foi atribuído a dois conceitos 3 e 4.
Os demais conceitos não foram atribuídos a nenhum curso e nenhum curso ficou Sem Conceito ( SC ) nesta combinação de categoria e região.
Os Centros Universitários da região Norte foram representados por dois cursos, aos quais foram atribuídos os conceitos 2 e 4.
Nenhum curso ficou Sem Conceito ( SC ) ou recebeu os demais conceitos nesta combinação de categoria e região.
As Faculdades participaram com oito cursos na região Norte que receberam os conceitos 2 ( seis cursos, conceito modal ) e 3 ( um curso ).
Um curso ficou Sem Conceito ( SC ).
Nenhuma das Faculdades da região obteve conceito 1, 4 ou 5.
Os CEFET/IFET participaram com cinco cursos.
O conceito modal nessa categoria foi 2, com três cursos, os demais cursos receberam os conceitos 3 e 4.
Na região Nordeste, as Universidades participaram com sete dos 47 cursos da Área de Tecnologia em Análise e Desenvolvimento de Sistemas oferecidos na região, dos quais um ficou Sem Conceito ( SC ), três cursos receberam, respectivamente, os conceitos 1, 2 e 4.
O conceito modal foi 3, atribuído a três cursos.
Nenhum curso oferecido por Universidades no Nordeste ficou com conceito 5.
Os Centros Universitários contaram com sete cursos participante na região Nordeste.
Esses obtiveram conceitos 2 ( um curso ), 3 ( quatro cursos, o conceito modal ) e conceito 4 ( dois cursos ).
Nenhum curso oferecido por Centros Universitários no Nordeste ficou Sem Conceito ( SC ) ou recebeu os conceitos 1 e 5.
As Faculdades foram representadas por 21 cursos na região Nordeste e o conceito modal foi 2, atribuído a nove cursos.
O conceito 1 foi atribuído a um curso, o conceito 3 a seis cursos, o conceito 4 a dois cursos e um curso recebeu conceito 5.
Na região Nordeste, dois cursos oferecido por Faculdades ficaram Sem Conceito ( SC ).
Os CEFET/IFET participaram com 12 cursos, os quais receberam conceitos 3 ( quatro cursos ), 4 ( cinco cursos, o conceito modal ) e 5 ( três cursos ).
Nenhum curso ficou Sem Conceito ( SC ), nem recebeu os conceitos 1 e 2.
Na região Sudeste, as Universidades concentraram 51 dos 189 cursos de Tecnologia em Análise e Desenvolvimento de Sistemas da região.
Entre os cursos em Universidades, nessa região, o conceito modal foi 3, atribuído a 23 cursos.
Os demais cursos receberam os conceitos 1 ( um curso ), 2 ( 21 cursos ), 4 ( quatro cursos ).
Dois cursos ficaram Sem Conceito ( SC ) e nenhum curso recebeu o conceito 5.
Os Centros Universitários participaram com 33 cursos na região Sudeste, dos quais 16 obtiveram o conceito modal 3.
O conceito 1 foi atribuído a dois cursos, o conceito 2, a sete cursos e o conceito 4, a seis cursos.
Nenhum curso recebeu conceito 5 e dois cursos ficaram Sem Conceito ( SC ) nessa combinação de organização acadêmica e região.
Já as Faculdades foram representadas por 77 cursos na região Sudeste, dos quais cinco ficaram Sem Conceito ( SC ).
A três cursos foram atribuídos, igualmente, os conceitos 1 e 5.
Aos demais cursos foram atribuídos os conceitos 2 ( 23 cursos ), 3 ( 29 cursos, o conceito modal ) e 4 ( 14 cursos ).
Os CEFET/IFET participaram com 28 cursos, os quais receberam conceitos 3 ( 13 cursos, o conceito modal ) e 4, ( 11 cursos ).
Os conceitos 1, 2 e 5 foram atribuídos a um curso cada e apenas um curso ficou Sem Conceito ( SC ).
Dos 91 cursos da região Sul, 28 eram de Universidades, para os quais o conceito modal foi 3, atribuído a 15 cursos.
Nesse tipo de organização, nenhum curso ficou Sem Conceito ( SC ) e nenhum curso recebeu conceito 1.
Os demais cursos receberam os conceitos 2 e 5 ( um curso, cada ) e 4 ( 11 cursos ).
Dos 15 cursos participantes de Centros Universitários da região Sul, nenhum recebeu os conceitos 1 e 5 e apenas um ficou Sem Conceito ( SC ).
O conceito modal, 4, foi atribuído a seis cursos.
Os conceitos restantes, 2 e 3, foram atribuídos a três e cinco cursos, respectivamente.
Quanto aos 30 cursos vinculados a Faculdades na região Sul, dez receberam o conceito modal 2.
O conceito 1 foi atribuído a um curso, o conceito 3, a nove cursos e o conceito 4 a seis cursos.
Nenhum curso recebeu o conceito 5 e quatro ficaram Sem Conceito ( SC ).
Nesta região 18 cursos foram avaliados em CEFET/IFET.
O conceito modal nessa categoria foi 4, atribuído a nove cursos.
Os demais cursos receberam conceito 5 ( três cursos ), 3 ( cinco cursos ), 2 ( um curso ) e nenhum curso ficou Sem Conceito ( SC ) ou com conceito 1 nessa região.
Na região Centro-Oeste, sete dos 42 cursos eram de Universidades.
Nesse tipo de organização, nenhum curso ficou Sem Conceito ( SC ) ou recebeu os conceitos 1 e 5.
Os conceitos 2 e 3 foram atribuídos a dois cursos cada e o conceito 4 ( o conceito modal ), a três cursos.
Os Centros Universitários da região Centro-Oeste contaram com dez cursos, que obtiveram os conceitos 3 ( quatro cursos, conceito modal ) e os conceitos 2 e 4, respectivamente, a três cursos cada.
Dos 14 cursos oferecidos em Faculdades na região Centro-Oeste, nove receberam o conceito modal 2, e quatro cursos ficaram com conceito 3.
O curso restante foi atribuído o conceito 4.
Nesta região 11 cursos foram avaliados em CEFET/IFET.
O conceito modal nessa categoria foi 3, atribuído a oito cursos.
Os demais cursos receberam conceito 2 ( três cursos ).
O Gráfico 5.3 apresenta a distribuição do Conceito Enade dos cursos de Licenciatura segundo a Organização Acadêmica da IES.
Os cursos em CEFET/IFET ( linha roxa ) apresentam uma poligonal mais à direita do que as demais, denotando uma melhor distribuição de Conceitos Enade.
No outro extremo, temos os cursos oferecidos em Faculdades ( linha vermelha ) com a poligonal mais à esquerda e, consequentemente, uma distribuição de Conceitos Enade com menores valores.
IES de Centros Universitários ( linha verde ) e em Universidades ( linha azul ) apresentam uma situação intermediária.

% matched lemmas: 
\end{textsample}
